% PAGES: 64-65  (printed 48-49)

To compute the first row of $U$, we multiply the first row of $L$ by the
column $k$ of $U$ and use the fact that $LU=A$ to obtain that
\begin{equation}
U(1,k) \;=\; A(1,k), \quad \forall\, k=1:n.
\end{equation}

In particular, for $k=1$ in (2.22), we find that
\begin{equation}
U(1,1) \;=\; A(1,1).
\end{equation}

To compute the first column of $L$, we multiply the row $k$ of $L$ by
the first column of $U$ and obtain that $L(k,1)U(1,1) = A(k,1)$ for all
$k=2:n$, and therefore
\begin{equation}
L(k,1) \;=\; \frac{A(k,1)}{U(1,1)} \;=\; \frac{A(k,1)}{A(1,1)}, \quad \forall\, k=1:n;
\end{equation}
see (2.23) for the last equality. Note that the fraction from (2.24) is
well defined: $A(1,1)\neq 0$ since $A(1,1)$ is the first leading
principal minor of $A$ and the matrix $A$ has an LU decomposition; cf.
Theorem 2.1 and Definition 2.1.

The other rows of $U$ and columns of $L$ are computed recursively, one
row of $U$ and one column of $L$ at a time, as follows:

To compute the second row of $U$ and the second column of $L$, write the
matrices $L$ and $U$ as
\begin{equation}
L \;=\; \begin{pmatrix} 1 & 0 \\ L(2:n,1) & L(2:n,2:n)\end{pmatrix},
\end{equation}
\begin{equation}
U \;=\; \begin{pmatrix} U(1,1) & U(1,2:n) \\ 0 & U(2:n,2:n)\end{pmatrix},
\end{equation}
where $L(2:n,1) = (L(j,1))_{j=2:n}$ is an $(n-1)\times 1$ column vector;
$U(1,2:n) = (U(1,k))_{k=2:n}$ is a $1\times(n-1)$ row vector;
$L(2:n,2:n) = (L(j,k))_{2\le j,k\le n}$ is the $(n-1)\times(n-1)$ lower
triangular matrix made of the entries from the rows $2, 3, \ldots, n$
and the columns $2, 3, \ldots, n$ of $L$; $U(2:n,2:n) =
(U(j,k))_{2\le j,k\le n}$ is the $(n-1)\times(n-1)$ upper triangular
matrix made of the entries from the rows $2, 3, \ldots, n$ and the
columns $2, 3, \ldots, n$ of $U$.

Similarly, write the matrix $A$ as
\begin{equation}
A \;=\; \begin{pmatrix} A(1,1) & A(1,2:n) \\ A(2:n,1) & A(2:n,2:n)\end{pmatrix},
\end{equation}
where
\[
\begin{array}{ll}
A(1,2:n) = (A(1,k))_{k=2:n} & \text{is an } 1\times(n-1) \text{ row vector;}\\
A(2:n,1) = (A(j,1))_{j=2:n} & \text{is an } (n-1)\times 1 \text{ column vector;}\\
A(2:n,2:n) = (A(j,k))_{2\le j,k\le n} & \text{is an } (n-1)\times(n-1) \text{ matrix.}
\end{array}
\]

From (2.25--2.27), it follows that $LU=A$ is equivalent to
\begin{equation}
\begin{pmatrix} 1 & 0 \\ L(2:n,1) & L(2:n,2:n)\end{pmatrix}
\begin{pmatrix} U(1,1) & U(1,2:n) \\ 0 & U(2:n,2:n)\end{pmatrix}
\end{equation}
\begin{equation}
\;=\; \begin{pmatrix} A(1,1) & A(1,2:n) \\ A(2:n,1) & A(2:n,2:n)\end{pmatrix}.
\end{equation}

By using block matrix multiplication to multiply the rows $2:n$ of $L$
by the columns $2:n$ of $U$, we obtain from (2.28--2.29)
that\footnote{Since $L(2:n,1)$ is an $(n-1)\times 1$ column vector and
$U(1,2:n)$ is a $1\times(n-1)$ row vector, it follows that the result of
the column vector -- row vector multiplication $L(2:n,1)U(1,2:n)$ is an
$(n-1)\times(n-1)$ matrix, see (1.6), which is the same size as the
matrices $A(2:n,2:n)$, $L(2:n,2:n)$, and $U(2:n,2:n)$. Thus, the matrix
dimensions in (2.30) are consistent.}
\begin{equation}
L(2:n,1)U(1,2:n) \;+\; L(2:n,2:n)U(2:n,2:n) \;=\; A(2:n,2:n),
\end{equation}
and therefore
\begin{equation}
L(2:n,2:n)U(2:n,2:n) \;=\; A(2:n,2:n) \;-\; L(2:n,1)U(1,2:n).
\end{equation}

From (2.31), we conclude that the matrices $L(2:n,2:n)$ and $U(2:n,2:n)$
are the $L$ and $U$ factors from the LU decomposition of the matrix
$A(2:n,2:n) - L(2:n,1)U(1,2:n)$. The first row of the matrix $U(2:n,2:n)$
(which coincides with the second row of $U$ without its first entry),
and the first column of the matrix $L(2:n,2:n)$ (which coincides with
the second column of $L$ without its first entry which is equal to $0$)
are computed from the matrix $A(2:n,2:n) - L(2:n,1)U(1,2:n)$ the same
way\footnote{Note that this step requires division by $U(2,2)$; see also
(2.24). The fact that $U(2,2) \neq 0$ and therefore the division can be
performed is a consequence of the fact that the leading principal minor
$\det(A(1:2,1:2))$ is nonzero; see also an exercise at the end of this
chapter.} as the first row of $U$ and the first column of $L$ were
computed from the matrix $A$.

Note that, in the implementation of the LU algorithm, we will call
computing the matrix $A(2:n,2:n) - L(2:n,1)U(1,2:n)$ ``updating'' the
matrix $A(2:n,2:n)$, and will use the notation
\begin{equation}
A(2:n,2:n) \;=\; A(2:n,2:n) \;-\; L(2:n,1)\,U(1,2:n),
\end{equation}
which means that we overwrite the old entries of $A(2:n,2:n)$ by the
entries of $A(2:n,2:n) - L(2:n,1)U(1,2:n)$.

The algorithm continues recursively until all the rows of $U$ and all
the columns of $L$ are computed.

\begin{examplex}
\textit{$4\times 4$ matrix}\par
% Continues onto page 50 (sec02_c2_4.tex): this examplex environment is
% closed there, after equation (2.35). Do not close it again.
For $n=4$, we obtain that
\[
A(1,2:n) \;=\; A(1,2:4) \;=\; \big(A(1,2)\ \ A(1,3)\ \ A(1,4)\big);
\]
\[
A(2:n,1) \;=\; A(2:4,1) \;=\; \begin{pmatrix} A(2,1)\\ A(3,1)\\ A(4,1)\end{pmatrix};
\]
\[
A(2:n,2:n) \;=\; A(2:4,2:4) \;=\;
\begin{pmatrix} A(2,2) & A(2,3) & A(2,4)\\ A(3,2) & A(3,3) & A(3,4)\\ A(4,2) & A(4,3) & A(4,4)\end{pmatrix};
\]
\[
L(2:n,1) \;=\; L(2:4,1) \;=\; \begin{pmatrix} L(2,1)\\ L(3,1)\\ L(4,1)\end{pmatrix};
\]
\[
L(2:n,2:n) \;=\; L(2:4,2:4) \;=\;
\begin{pmatrix} 1 & 0 & 0\\ L(3,2) & 1 & 0\\ L(4,2) & L(4,3) & 1\end{pmatrix};
\]
